% =====================================================================================================
\section{Overview}\label{sec:overview}

This section states the results in words and says where each one is. The data are $n$ observations of a response and $p$
predictors, with $p$ possibly much larger than $n$. The response follows a sparse Gaussian linear model, and the method
is a sequential Monte Carlo sampler for the posterior over models under a fixed prior, combined with a local search.

\paragraph{Does the posterior pick the right model? (Section~\ref{sec:selection})} Yes. For fixed designs, the posterior
mass outside the true model is at most a constant times $\rn$ on average, uniformly over the designs that meet the
conditions (Theorems~\ref{thm:post-fixed} and~\ref{thm:post-margin}); the highest posterior model and the median
probability model equal the true model with probability tending to one. The same holds for random Gaussian designs
(Theorem~\ref{thm:post-random}). No design allows a rate faster than $n^{-1/2}$, and on two families of designs the rate
$\rn$ is attained (Theorem~\ref{thm:sharp}). When some true predictors are missing from the data, the posterior selects
the predictors used by the best linear predictor built from the others (Theorem~\ref{thm:omitted}).

\paragraph{Is the posterior right about the distribution? (Section~\ref{sec:distribution})} Yes. The posterior contracts
around the true conditional law of the response given the predictors, at the rate $\rn$, in the Hellinger distance of
\citet{SongLiang2015}, with no beta-min condition (Theorem~\ref{thm:dist}); at the parametric rate when the true
predictors are detectable (Theorem~\ref{thm:dist-param}); and the same holds for the list of models that a completed run
returns (Corollary~\ref{cor:dist-output}).

\paragraph{Does the method, as coded, deliver it? (Section~\ref{sec:delivery})} Yes, under conditions on how the signal
looks block by block. The \textbf{main theorem}, Theorem~\ref{thm:reported-random}, says that on one data event of
probability tending to one, every run that produces output reports exactly the true model, whatever its seed. On the way,
the method's start law already concentrates on the true model (Theorem~\ref{thm:start}). A certified small model rarely
misses a true predictor, under much weaker conditions (Theorems~\ref{thm:screen} and~\ref{thm:screen-margin}). Every
reported estimate is consistent, and the evidence estimate is unbiased (Theorem~\ref{thm:outputs},
Remark~\ref{rem:unbiased}).

\paragraph{What carries over to the parallel method? (Section~\ref{sec:parallel})} The pooling of the four islands is
covered; the rest is empirical or open, and we say so.

\paragraph{What ``asymptotic'' means here.} Every result is a statement as $n\to\infty$, with $p$ allowed to grow
polynomially in $n$ and to exceed $n$. At moderate $n$ and large $p$ the posterior is not close to a point mass: it is
spread over the true model and its many one-predictor supersets, each individually light (Remark~\ref{rem:moderate}). The
theorems describe the limit of this picture.

\paragraph{How to read a result.} Each result has the same layout: a sentence on the question it answers; a paragraph
\emph{In words}; the new assumptions it needs, each with a line on what it rules out; the formal statement; a shaded
status note; and remarks with fixed titles. Assumptions used by two or more theorems have short names, such as \as{GL};
assumptions used by one theorem are numbered with it, such as Assumption~\ref{asm:margin}. Appendix~\ref{app:notation}
lists every symbol and every assumption with the theorems that use it. ``The method'' always means the method as coded,
under Convention~\as{C1}.

\paragraph{Status notes.} Each shaded note gives the verification status of the result: ``proved and independently
checked'', ``proved and checked, conditional on'' named inputs, or ``in progress''. Nothing in progress is called a
theorem. The notes are for internal use and are removed for submission; Appendix~\ref{app:status} gives the records.

\begin{table}[ht]
\centering\small
\caption{The results at a glance.}\label{tab:overview}
\begin{tabularx}{\linewidth}{@{}l>{\raggedright\arraybackslash}X l l l@{}}
\toprule
Result & In words & Design & Rate & Status \\
\midrule
Thm~\ref{thm:post-fixed} & posterior concentrates on the true model & fixed & $\rn$ & checked$^{a,e}$ \\
Thm~\ref{thm:post-margin} & the same under an identifiability margin & fixed & $\rn$ & conditional$^{b,e}$ \\
Thm~\ref{thm:post-random} & the same for Gaussian rows & random & $\rn$ & checked$^{e}$ \\
Thm~\ref{thm:sharp} & the rate cannot be improved & both & $n^{-1/2}$; $\rn$ & checked$^{e}$ \\
Thm~\ref{thm:omitted} & omitted variables: support of the best candidate predictor & random & $\rn$ & conditional$^{b,e}$ \\
Thm~\ref{thm:dist} & posterior recovers the conditional law & random & $\rn$ & checked \\
Thm~\ref{thm:dist-param} & parametric rate for detectable signals & random & $M_n/\sqrt n$ & checked$^{c}$ \\
Cor~\ref{cor:dist-output} & the same for the list a completed run returns & random & $\rn$ & checked$^{d}$ \\
Thm~\ref{thm:start} & the start law concentrates on the true model & random & $n^{-c}$; $\rn$ & checked \\
Thm~\ref{thm:reported-random} & \textbf{every run that produces output reports the true model} & random & --- & checked$^{e}$ \\
Thm~\ref{thm:reported-fixed} & the same, fixed design & fixed & --- & checked$^{e}$ \\
Thms~\ref{thm:screen}, \ref{thm:screen-margin} & a certified small model rarely misses a true predictor & fixed & --- & checked \\
Thm~\ref{thm:outputs} & every reported estimate is consistent & random & $\rho_n$ & checked$^{e}$ \\
Rem~\ref{rem:unbiased} & the pooled evidence estimate is unbiased & any & exact & see note$^{f}$ \\
Sec~\ref{sec:parallel} & parallel engine; many islands & --- & --- & in progress / open \\
\bottomrule
\end{tabularx}
\par\smallskip\raggedright\footnotesize
$^{a}$ In-probability form; the expectation bound is conditional on two supporting lemmas not yet independently checked.
$^{b}$ Conditional on those two lemmas.
$^{c}$ Its one input, Theorem~\ref{thm:post-fixed} in probability form, is checked.
$^{d}$ Under exact arithmetic.
$^{e}$ The short bridging argument that puts the statement in reader form awaits its final independent check.
$^{f}$ Per island checked; the pooled form is not yet independently checked.
\end{table}

\paragraph{Correspondence with \citet{SongLiang2015}.} The table points to where each comparison is made; it adds no claim.

\begin{table}[ht]
\centering\small
\caption{Where each result is compared with \citet{SongLiang2015}.}\label{tab:SL}
\begin{tabular}{@{}lll@{}}
\toprule
This paper & Song and Liang (2015) & Comparison \\
\midrule
Theorems~\ref{thm:post-fixed}, \ref{thm:post-margin} & Theorem 3 (sure screening) & Remark~\ref{rem:SL-fixed} \\
Theorem~\ref{thm:post-random} & random-design setting & Remark~\ref{rem:post-random}(iii) \\
Theorem~\ref{thm:omitted} & Theorem 4, Corollary 2 & Remark~\ref{rem:SL-omitted} \\
Theorems~\ref{thm:dist}, \ref{thm:dist-param}, Corollary~\ref{cor:dist-output} & Theorems 1--2 & Section~\ref{sec:SL-dist} \\
Theorem~\ref{thm:reported-random} & the SAMC sampler they run & Remark~\ref{rem:SL-reported} \\
\bottomrule
\end{tabular}
\end{table}
